\begin{tikzpicture}[
  cathead/.style={
    draw, thick, align=center, font=\scriptsize\bfseries,
    inner sep=6pt, text width=40mm, minimum height=14mm,
  },
  catstep/.style={
    draw, rounded corners=2pt, align=center, font=\scriptsize,
    inner sep=5pt, text width=38mm, minimum height=16mm,
  },
  catgap/.style={
    draw, dashed, rounded corners=2pt, align=center, font=\scriptsize,
    inner sep=5pt, text width=38mm, minimum height=16mm,
  },
  catnote/.style={
    rectangle split,
    rectangle split parts=3,
    rectangle split part align=left,
    draw, dashed, rounded corners=2pt,
    font=\scriptsize, align=left,
    inner xsep=6pt, inner ysep=4pt,
    text width=120mm,
  },
]
  \node[cathead] (head) {CAT14\\translation loss};

  \node[catgap, below=16mm of head] (bridge)
    {Taught bridge?\\[2pt]{\tiny often missing}};
  \node[catstep, left=12mm of bridge] (src)
    {Source family\\[2pt]{\tiny iStar / BPMN / RE}};
  \node[catstep, right=12mm of bridge] (tgt)
    {Target family\\[2pt]{\tiny UML classes}};

  \node[catnote, below=12mm of bridge] (loss) {%
    \textbf{Often survives}\\names, requirements
    \nodepart{two}
    \textbf{Often drops}\\goals, process structure, and which question the model answers (S6, S10)
    \nodepart{three}
    \textbf{Bound}\\S7 is a lighter syntax-at-switch variant.%
  };

  \draw[thick] (head.south) -- (bridge.north);
  \draw[arr] (src) -- (bridge);
  \draw[arr] (bridge) -- (tgt);
\end{tikzpicture}
